% LaTeX Canvas: Chapter Outline
\chapter{Version Control}
\label{ch:git-github}

\section*{Purpose}
Version control is the safety system for technical work. Git records a history of changes so you can recover, compare, and collaborate without overwriting each other. GitHub adds shared hosting, review workflows, and discussion tools. This chapter teaches novice-friendly Git fundamentals and the everyday GitHub workflows students need: pushing and pulling, branches and merges, pull requests, forks, and review/commenting practices.

\section*{Learning objectives}
By the end of this chapter, a reader should be able to:
\begin{enumerate}
\item Explain what version control is and why Git is ``distributed.''
\item Create or clone a repository and make clean commits.
\item Understand the working tree, staging area, and commit history.
\item Use branches for features/fixes and merge changes safely.
\item Work with remotes (\texttt{origin}, \texttt{upstream}) and synchronize using fetch/pull/push.
\item Use GitHub to open a pull request, respond to feedback, and merge.
\item Resolve common merge conflicts with a disciplined playbook.
\item Use forks for contributing to projects you do not control.
\item Write good commit messages and PR descriptions; comment constructively in reviews.
\end{enumerate}

\section*{Running theme: make changes small, reviewable, and reversible}
A healthy Git workflow favors small commits, small pull requests, clear messages, and frequent synchronization.

\section{Mental models and vocabulary}
\subsection{Repository, snapshot, history}
\begin{itemize}
\item A \textbf{repository} is your project folder plus a hidden \texttt{.git} directory.
\item A \textbf{commit} is a snapshot of tracked files with metadata (author, time, message).
\item Git history is a graph, not a single line.
\end{itemize}

\subsection{The three zones: working tree, staging area, commits}
\begin{itemize}
\item \textbf{Working tree:} your edited files.
\item \textbf{Staging area (index):} the exact changes you plan to commit.
\item \textbf{Commit:} saved snapshot in history.
\end{itemize}

\subsection{Branching as pointers}
\begin{itemize}
\item A \textbf{branch} is a movable pointer to a commit.
\item A branch is cheap; use them.
\end{itemize}

\subsection{Remote repositories}
\begin{itemize}
\item \textbf{Local repo:} on your computer.
\item \textbf{Remote repo:} hosted on GitHub.
\item \textbf{origin:} the default remote name.
\item \textbf{upstream:} the source repository for forks.
\end{itemize}

\subsection{Pull request (PR)}
\begin{itemize}
\item A PR is a proposal to merge changes, with discussion and review.
\item PRs are where teams make decisions visible.
\end{itemize}

\section{Why version control matters for student data work}
\subsection{Use cases}
\begin{itemize}
\item Recover from mistakes (accidental edits, deletions).
\item Understand what changed and when.
\item Collaborate without overwriting.
\item Create an audit trail for decisions.
\end{itemize}

\subsection{What to version and what not to version}
\begin{itemize}
\item Version: code, notebooks (with discipline), docs, small configuration files.
\item Avoid committing: secrets, credentials, large raw datasets (unless instructed), generated artifacts.
\end{itemize}

\section{Getting started: create or clone}
\subsection{Initialize a repository (local)}
\begin{itemize}
\item When you start a project from scratch.
\item Add a \texttt{README.md} and a reasonable \texttt{.gitignore} early.
\end{itemize}

\subsection{Clone an existing repository}
\begin{itemize}
\item When the course/lab provides starter code.
\item Understand that cloning includes full history.
\end{itemize}

\subsection{First-time setup (identity and defaults)}
\begin{itemize}
\item Set your name/email.
\item Choose a default branch name consistent with the course/lab.
\end{itemize}

\section{Daily Git workflow (the minimum viable loop)}
\subsection{Status and diff first}
\begin{itemize}
\item Check what changed before committing.
\item Review diffs so commits are intentional.
\end{itemize}

\subsection{Stage intentionally}
\begin{itemize}
\item Stage the smallest coherent change.
\item Avoid committing unrelated changes together.
\end{itemize}

\subsection{Commit with meaning}
\begin{itemize}
\item Commit messages: short summary + optional details.
\item Prefer imperative mood: ``Add data intake checks''.
\end{itemize}

\subsection{Sync with the remote}
\begin{itemize}
\item Push your work so it is backed up and shareable.
\item Pull (or fetch+merge) so you do not drift from teammates.
\end{itemize}

\section{Branches and merges (collaboration without chaos)}
\subsection{Why branches}
\begin{itemize}
\item Keep \texttt{main} stable.
\item Work on features and fixes in isolation.
\end{itemize}

\subsection{A standard branch workflow}
\begin{enumerate}
\item Update \texttt{main}.
\item Create a new branch with a descriptive name.
\item Make commits on that branch.
\item Open a PR.
\item Merge after review.
\end{enumerate}

\subsection{Merge vs rebase (novice policy)}
\begin{itemize}
\item \textbf{Merge} preserves history of parallel work.
\item \textbf{Rebase} rewrites commits to create a linear story.
\item Student default: use merges unless your instructor/team specifies a rebase policy.
\end{itemize}

\subsection{Undo safely (intro)}
\begin{itemize}
\item Undo a working-tree change vs undo a commit are different.
\item Prefer ``make a new commit that fixes it'' over rewriting shared history.
\end{itemize}

\section{Remotes: fetch, pull, push}
\subsection{What these commands do conceptually}
\begin{itemize}
\item \texttt{fetch}: download remote history without changing your current branch.
\item \texttt{pull}: fetch + merge (or fetch + rebase).
\item \texttt{push}: upload your commits to the remote.
\end{itemize}

\subsection{The drift problem}
\begin{itemize}
\item If multiple people push, your local copy can become stale.
\item Best practice: fetch/pull frequently; push small, complete units.
\end{itemize}

\section{GitHub fundamentals}
\subsection{Repository anatomy on GitHub}
\begin{itemize}
\item Code view, issues, pull requests, actions, releases.
\item The README is the front door.
\end{itemize}

\subsection{Pull requests: the core collaboration primitive}
\begin{itemize}
\item What a PR contains: commits, file diffs, discussion.
\item Draft PRs for early feedback.
\item Keep PRs small and focused.
\end{itemize}

\subsection{PR merge strategies (team policy)}
\begin{itemize}
\item Merge commit vs squash-and-merge vs rebase-and-merge.
\item Student guidance: pick one strategy per repository and be consistent.
\end{itemize}

\subsection{Reviews and commenting best practices}
\begin{itemize}
\item Comment on code with context: what, why, and suggested change.
\item Prefer actionable, specific feedback.
\item Separate `nit'' (minor style) from `blocker'' (correctness).
\item When you respond, acknowledge and summarize what you changed.
\end{itemize}

\subsection{Linking work: issues $\leftrightarrow$ PRs}
\begin{itemize}
\item Use issues to track tasks and decisions.
\item Reference issues in PR descriptions and commit messages.
\end{itemize}

\section{Forking workflow (contributing without direct write access)}
\subsection{Fork vs clone}
\begin{itemize}
\item Fork creates your copy on GitHub.
\item Clone downloads a repository to your computer.
\end{itemize}

\subsection{The standard fork contribution flow}
\begin{enumerate}
\item Fork the repository.
\item Clone your fork.
\item Add \texttt{upstream} remote pointing to the original repo.
\item Create a feature branch.
\item Push to your fork.
\item Open a PR from your fork to upstream.
\end{enumerate}

\subsection{Keeping your fork up to date}
\begin{itemize}
\item Fetch from upstream.
\item Merge (or rebase) into your branch.
\item Resolve conflicts early rather than letting them accumulate.
\end{itemize}

\section{Merge conflicts: a disciplined resolution playbook}
\subsection{What a conflict is}
\begin{itemize}
\item Git cannot automatically reconcile changes to the same lines/regions.
\item Conflicts are normal in collaboration.
\end{itemize}

\subsection{The safe sequence}
\begin{enumerate}
\item Stop and read the error.
\item Identify which files are conflicted.
\item Open the files and locate conflict markers.
\item Decide the correct combined content (not ``pick one side" by default).
\item Stage resolved files.
\item Complete the merge.
\item Run a quick smoke test.
\end{enumerate}

\subsection{When to ask for help}
\begin{itemize}
\item If you cannot explain what each side of a conflict represents.
\item If conflicts involve data files or notebooks that are difficult to merge.
\end{itemize}

\section{Special considerations for data science projects}
\subsection{\texttt{.gitignore} hygiene}
\begin{itemize}
\item Ignore environment folders (\texttt{.venv/}, conda env dirs).
\item Ignore caches and generated outputs (\texttt{__pycache__}, \texttt{.ipynb_checkpoints/}).
\item Ignore secrets files (\texttt{.env}).
\end{itemize}

\subsection{Notebooks in Git}
\begin{itemize}
\item Notebooks are JSON; diffs can be noisy.
\item Practices: clear unnecessary outputs, keep cells ordered, consider paired text formats.
\end{itemize}

\subsection{Large files and datasets}
\begin{itemize}
\item Avoid committing large raw data unless instructed.
\item Prefer storing raw data outside the repo and documenting retrieval.
\end{itemize}

\subsection{Releases and tags (optional)}
\begin{itemize}
\item Tag a version that corresponds to a submission or report.
\item Helps you reproduce exactly what was delivered.
\end{itemize}

\section{Worked examples (outline)}
\subsection{Example 1: The daily loop}
\begin{itemize}
\item Edit file, check status/diff, stage, commit, push.
\end{itemize}

\subsection{Example 2: Branch + PR workflow}
\begin{itemize}
\item Create feature branch, commit in small steps, open PR, respond to review, merge.
\end{itemize}

\subsection{Example 3: Resolve a merge conflict}
\begin{itemize}
\item Create an intentional conflict, resolve using the playbook, run a smoke test.
\end{itemize}

\subsection{Example 4: Fork contribution}
\begin{itemize}
\item Fork, clone, add upstream, sync, open PR to upstream.
\end{itemize}

\section{Templates}
\subsection{Template A: Commit message}
\begin{verbatim} <Short summary in imperative mood>

Why:

* ...

What changed:

* ...
  \end{verbatim}

\subsection{Template B: Pull request description}
\begin{verbatim}

## Summary

## Why

## What changed

*

## How to test

*

## Screenshots/outputs (if relevant)

## Related issues

Closes #
\end{verbatim}

\subsection{Template C: Review comment taxonomy}
\begin{verbatim}
Blocker: correctness, security, reproducibility
Suggestion: improvement, alternative approach
Question: clarify intent, edge cases
Nit: formatting, naming (non-blocking)
\end{verbatim}

\section{Exercises}
\begin{enumerate}
\item Initialize a repo, add a README and \texttt{.gitignore}, and make three meaningful commits.
\item Create a branch, make changes, push it, and open a PR.
\item Review a classmate's PR using the comment taxonomy; request one change and approve after revision.
\item Create a merge conflict deliberately and resolve it using the playbook.
\item Fork a repository, add an upstream remote, sync it, and open a PR from your fork.
\end{enumerate}

\section{One-page checklist}
\begin{itemize}
\item I understand working tree vs staging vs commits.
\item I can create branches and keep \texttt{main} stable.
\item I can fetch/pull/push and explain what each does.
\item I can open, review, and merge pull requests.
\item I can use forks and keep them synchronized with upstream.
\item I can resolve merge conflicts methodically.
\item I follow hygiene: \texttt{.gitignore}, no secrets, controlled notebook diffs.
\end{itemize}

\appendix
\section{Quick reference: common commands (student set)}
\begin{verbatim}

# create/clone

git init
git clone <url>

# inspect

git status
git diff
git log --oneline --graph --decorate

# stage/commit

git add <file>
git add -p
git commit -m "..."

# branches

git branch
git switch -c <name>
git switch <name>

# remotes

git remote -v
git fetch
git pull
git push -u origin <branch>

# merge/conflicts

git merge <branch>
git merge --abort
\end{verbatim}
